\documentclass[12pt,a4paper]{article}

\usepackage{amsmath,amssymb,amsthm}
\usepackage{mathtools}
\usepackage{geometry}
\usepackage{doi}
\geometry{margin=2.5cm}
\usepackage{hyperref}
\usepackage{cleveref}
\usepackage{enumitem}

% -------------------------------------------------------------------
% Theorem environments
% -------------------------------------------------------------------
\newtheorem{theorem}{Theorem}[section]
\newtheorem{proposition}[theorem]{Proposition}
\newtheorem{lemma}[theorem]{Lemma}
\newtheorem{corollary}[theorem]{Corollary}
\theoremstyle{definition}
\newtheorem{definition}[theorem]{Definition}
\newtheorem{remark}[theorem]{Remark}

% -------------------------------------------------------------------
% Macros
% -------------------------------------------------------------------
\newcommand{\Spi}{S_{\Pi}}
\newcommand{\Ag}{A_g}
\newcommand{\hk}{K}
\newcommand{\tr}{\operatorname{Tr}}
\newcommand{\vol}{\operatorname{vol}}
\newcommand{\dd}{\mathrm{d}}
\newcommand{\half}{\tfrac{1}{2}}

\begin{document}

  \title{Local Spectral Thermodynamics and the Einstein Equation as a Spectral Equilibrium Condition}

  \author{J\'{e}r\^{o}me Beau\\
  \small Independent Researcher, France\\
  \small \texttt{jerome.beau@cosmochrony.org}}

  \maketitle

  \begin{abstract}
    We develop a local thermodynamic interpretation of the projective spectral entropy
    functional $\Spi[g] = \frac{1}{2}\log\det' A_g$, extending the covariant framework
    established in a companion paper.
    The diagonal heat kernel $K(x,x;t)$ defines a local spectral energy density
    $u(x;t) = -\partial_t \log K(x,x;t)$, whose Seeley--DeWitt expansion yields
    a local spectral multiplier field $\beta(x)^{-1} = \beta_0^{-1} - \frac{1}{6}R(x) + O(\nabla^2 R)$.
    Here $\beta_0^{-1} = 2/t_*$ is a universal UV contribution fixed by the spectral
    admissibility scale $t_*$, while the curvature correction is geometric and covariant.

    We then formulate a local spectral first law as a compatibility theorem:
    in the infrared-dominant regime, the metric variation of $\Spi$ and the metric
    variation of the projected matter functional are related through the local multiplier
    field, yielding a constraint $\delta\Spi = \int \beta(x)^{-1}\,\delta E_\Pi(x)$.
    This is not a definition but a structural consequence of the Seeley--DeWitt
    hierarchy and the induced-gravity normalization established in the companion paper.

    We then impose a spectral equilibrium principle: admissible geometries extremize
    $\Spi$ at fixed projected matter content.
    The resulting Euler--Lagrange equation is Einstein's field equation
    $G_{\mu\nu} = 8\pi G\, T_{\mu\nu}$, with the factor $8\pi G$ arising purely from
    the normalization conventions of the geometric and matter sides.
    The derivation requires no horizon structure, no Rindler wedge, no Raychaudhuri
    theorem, and no assumption of local Lorentz invariance at the Planck scale.
    Einstein's equation emerges as a spectral extremum condition within the
    effective infrared geometry.
  \end{abstract}

  \noindent\textbf{Keywords:} spectral entropy, heat kernel, local temperature,
  induced gravity, Einstein equation, spectral equilibrium, Seeley--DeWitt

  \clearpage
  \tableofcontents
  \newpage

  \section{Introduction}
    \label{sec:intro}

    In a companion paper~\cite{Beau2026h}, we established that the projective spectral
    entropy functional
    \begin{equation}
      \Spi[g] \equiv \half\log\det' \Ag ,
      \label{eq:Spi-def}
    \end{equation}
    defined by a minimal Laplace-type operator $\Ag = -\nabla_g^2$ on a four-dimensional
    Riemannian manifold $(M,g)$, admits a well-defined renormalized metric variation.
    In the infrared-dominant regime where all dimensionless curvature invariants satisfy
    $R\ell_\chi^2 \ll 1$, the renormalized spectral stress tensor decomposes as
    \begin{equation}
      \frac{\delta\Spi^{\mathrm ren}}{\delta g^{\mu\nu}}
      = c_{\mathrm EH}({\mu})\,G_{\mu\nu} + c_\Lambda(\mu)\,g_{\mu\nu}
      + \beta_{\mathrm loc}\,B_{\mu\nu} + T^{\mathrm non\text{-}loc}_{\mu\nu} + A_{\mu\nu},
      \label{eq:decomp-paper1}
    \end{equation}
    with $c_{\mathrm EH} = (16\pi G_N)^{-1}$ identified via induced-gravity renormalization;
    higher-derivative terms are suppressed by powers of $\ell_\chi^2$.

    The goal of the present paper is to develop the local thermodynamic interpretation
    of this structure.
    We show that the diagonal heat kernel $\hk(x,x;t)$ defines a natural local spectral
    energy density $u(x;t)$, whose curvature expansion produces a local multiplier field
    $\beta(x)^{-1}$ with a universal part and a geometric correction proportional to
    the scalar curvature $R(x)$.
    This field plays the role of a local inverse temperature, defined purely from spectral
    data, without thermodynamic input.

    We then show that the compatibility between the geometric and matter sides of the
    metric variation is precisely the condition
    \begin{equation}
      \delta\Spi = \int_\Sigma \beta(x)^{-1}\,\delta E_\Pi(x)\,\dd^3 x ,
      \label{eq:first-law-intro}
    \end{equation}
    which we call the \emph{local spectral first law}.
    This is not a postulate but a structural consequence of the Seeley--DeWitt
    hierarchy.

    Finally, we impose a spectral equilibrium principle: admissible geometries are
    those for which $\Spi$ is stationary at fixed projected matter content.
    The Euler--Lagrange equation of this constrained extremum is Einstein's field
    equation, obtained without any reference to local horizons, Rindler observers,
    or causal structure.

    \medskip
    \noindent\textbf{Relation to Jacobson's derivation.}
    Jacobson~\cite{Jacobson1995} derived Einstein's equation by imposing the Clausius
    relation $\delta S = \delta Q/T$ on all local Rindler horizons, using the
    Bekenstein--Hawking entropy proportional to area.
    The present derivation is structurally parallel but differs in three essential ways:
    (i) the entropy functional is $\Spi$, not an area functional;
    (ii) the temperature field $\beta(x)^{-1}$ is defined from spectral data, not from
    Rindler kinematics;
    (iii) the field equation arises from a variational extremum, not from a horizon
    equilibrium condition.
    The relation to Jacobson's derivation is discussed further in
    Section~\ref{sec:relation-jacobson}.

    \medskip
    \noindent\textbf{Structure of the paper.}
    Section~\ref{sec:spectral-locality} develops the local spectral energy density and
    the multiplier field $\beta(x)$.
    Section~\ref{sec:first-law} establishes the local spectral first law as a
    compatibility theorem.
    Section~\ref{sec:equilibrium} formulates the spectral equilibrium principle and
    derives the Einstein equation.
    Section~\ref{sec:normalization} verifies the normalization of the factor $8\pi G$.
    Section~\ref{sec:relation-jacobson} discusses the relation to previous thermodynamic
    derivations.
    Section~\ref{sec:discussion} collects limitations and outlook.

  \section{Spectral locality and the temperature multiplier field}
    \label{sec:spectral-locality}

    \subsection{Local spectral energy density}
      \label{subsec:local-density}

      Let $\Ag$ be a positive Laplace-type operator on $(M,g)$.
      Denote by $\hk(x,y;t)$ its heat kernel, satisfying
      \begin{equation}
        (\partial_t + \Ag)\hk(x,y;t) = 0,
        \qquad
        \hk(x,y;0^+) = \delta_g(x,y).
        \label{eq:heat-eq}
      \end{equation}
      The diagonal $u_0(x;t) \equiv \hk(x,x;t)$ is a local scalar density encoding
      the spectral weight at point $x$ and diffusion scale $t$.

      \begin{definition}[Local spectral energy density]
        \label{def:u}
        The \emph{local spectral energy density} at point $x$ and scale $t$ is
        \begin{equation}
          u(x;t) \equiv -\partial_t \log \hk(x,x;t).
          \label{eq:u-def}
        \end{equation}
      \end{definition}

      This definition is minimal: it depends only on local diagonal heat-kernel data.
      The sign convention ensures $u(x;t) > 0$ in the short-time regime, consistent with
      an energy interpretation.
      The expression is also the natural local analogue of the canonical relation
      $E = -\partial_\beta \log Z$ in statistical mechanics, once $t$ is identified with
      an inverse spectral resolution parameter.

    \subsection{Seeley--DeWitt expansion of \texorpdfstring{$u(x;t)$}{u(x;t)}}
      \label{subsec:SDW-expansion}

      For the minimal scalar Laplacian $\Ag = -\nabla_g^2$ in dimension $d=4$, the
      short-time expansion of the diagonal heat kernel reads~\cite{Vassilevich2003,
    Gilkey1995}
      \begin{equation}
        \hk(x,x;t) \sim \frac{1}{(4\pi t)^2}
        \bigl(a_0(x) + a_2(x)\,t + a_4(x)\,t^2 + \cdots\bigr)
        \quad\text{as }t\to 0^+,
        \label{eq:hk-expansion}
      \end{equation}
      with Seeley--DeWitt coefficients
      \begin{align}
        a_0(x) &= 1, \label{eq:a0}\\
        a_2(x) &= \tfrac{1}{6}R(x), \label{eq:a2}\\
        a_4(x) &= \frac{1}{(4\pi)^2\cdot 360}
        \bigl(R_{\mu\nu\rho\sigma}R^{\mu\nu\rho\sigma}
        - R_{\mu\nu}R^{\mu\nu} + \tfrac{5}{3}R^2\bigr). \label{eq:a4}
      \end{align}

      Taking the logarithm of~\eqref{eq:hk-expansion} and differentiating with respect
      to $t$ yields:
      \begin{equation}
        \partial_t \log \hk(x,x;t)
        = -\frac{2}{t}
        + \partial_t \log\bigl(a_0 + a_2\,t + a_4\,t^2 + \cdots\bigr).
        \label{eq:dlog-hk}
      \end{equation}
      For the minimal Laplacian, $a_0 = 1$ and $a_2 = \frac{1}{6}R$, so
      \begin{equation}
        \partial_t \log\bigl(1 + \tfrac{1}{6}R\,t + O(t^2)\bigr)
        = \tfrac{1}{6}R + O(t).
        \label{eq:log-expand}
      \end{equation}
      Therefore, by Definition~\ref{def:u}:

      \begin{proposition}[Curvature expansion of $u$]
        \label{prop:u-expansion}
        For the minimal scalar Laplacian in dimension four,
        \begin{equation}
          u(x;t) = \frac{2}{t} - \frac{1}{6}R(x) + O(t).
          \label{eq:u-expansion}
        \end{equation}
        The leading term is universal and UV-divergent; the first geometric correction is
        proportional to the scalar curvature with coefficient fixed by $a_2$.
      \end{proposition}

    \subsection{The local multiplier field \texorpdfstring{$\beta(x)$}{beta(x)}}
      \label{subsec:beta-field}

      We now evaluate $u(x;t)$ at a fixed spectral scale $t = t_*$, defined as the
      admissibility scale introduced in~\cite{Beau2026h,Beau2026a} as the infrared
      threshold below which the spectral description is consistent.

      \begin{definition}[Local spectral multiplier field]
        \label{def:beta}
        The \emph{local spectral multiplier field} $\beta(x)^{-1}$ is defined by
        \begin{equation}
          \beta(x)^{-1} \equiv u(x;t_*) = \frac{2}{t_*} - \frac{1}{6}R(x) + O(t_*).
          \label{eq:beta-def}
        \end{equation}
      \end{definition}

      \begin{remark}[Sign convention and decomposition]
        \label{rem:sign}
        Equation~\eqref{eq:beta-def} may appear to give a \emph{negative} curvature correction,
        whereas the literature on local temperature often writes a \emph{positive} geometric
        correction.
        The sign is consistent and arises from the definition $u = -\partial_t \log K$
        together with the expansion $\log K \sim -2\log t + \frac{1}{6}Rt + \cdots$,
        which gives $\partial_t \log K \sim -2/t + \frac{1}{6}R$,
        and hence $u \sim 2/t - \frac{1}{6}R$.

        When we write the decomposition with $\beta_0^{-1} \equiv 2/t_*$, we have
        \begin{equation}
          \beta(x)^{-1} = \beta_0^{-1} - \frac{1}{6}R(x) + O(t_*).
          \label{eq:beta-decomp}
        \end{equation}
        On a positively curved manifold ($R > 0$), the local multiplier is \emph{smaller}
        than the flat reference $\beta_0^{-1}$: curvature reduces the effective spectral
        temperature relative to flat space.
        This sign is physically consistent with the Seeley--DeWitt interpretation:
        positive curvature lowers the local spectral density, hence lowers the effective
        temperature.
      \end{remark}

      \begin{remark}[Status of $t_*$ and $\beta_0^{-1}$]
        \label{rem:tstar}
        The scale $t_*$ is a spectral admissibility parameter, not a physical temperature.
        Its value is fixed by the spectral resolution of the relational system, as discussed
        in~\cite{Beau2026a}.
        The universal part $\beta_0^{-1} = 2/t_*$ is a UV quantity reflecting the spectral
        density of states at the cutoff.
        It is scheme-dependent and carries no intrinsic thermodynamic meaning at this stage.
        Only the geometric correction $\frac{1}{6}R(x)$ is covariant and physically
        significant: it encodes local curvature as a deviation from the flat-space multiplier.
      \end{remark}

      \begin{remark}[Decomposition structure]
        \label{rem:decomp}
        The decomposition~\eqref{eq:beta-decomp} has the form
        \[
          \beta(x)^{-1}
          = \underbrace{\beta_0^{-1}}_{\text{universal (UV)}}
          - \underbrace{\tfrac{1}{6}R(x)}_{\text{geometric correction}}.
        \]
        The geometric correction carries a negative sign: positive curvature reduces the
        local spectral multiplier relative to the flat-space reference $\beta_0^{-1}$.
        The correction is entirely controlled by the $a_2$ Seeley--DeWitt coefficient,
        which is precisely the coefficient that generates the Einstein--Hilbert counterterm
        via pole subtraction in the companion paper.
        The two papers are therefore connected at the level of the same spectral coefficient.
      \end{remark}

  \section{Local spectral first law as a compatibility theorem}
    \label{sec:first-law}

    \subsection{Projected matter energy}
      \label{subsec:projected-matter}

      We define the projected matter energy on a spatial slice $\Sigma$ independently
      of the spectral multiplier field.
      Let $W_\Pi[g,\psi]$ be the effective matter action associated with the admissible
      projected sector.
      Its metric variation defines the projected stress tensor $T^{(\Pi)}_{\mu\nu}$ via
      the standard covariant identity
      \begin{equation}
        \delta W_\Pi[g,\psi]
        = \half\int_M \dd^4x\,\sqrt{g}\;T^{(\Pi)}_{\mu\nu}\,\delta g^{\mu\nu}.
        \label{eq:W-variation}
      \end{equation}
      On a chosen spatial slice $\Sigma$ with future-pointing unit normal $n^\mu$ and
      local Killing vector $\xi^\mu$, the projected matter energy is defined by
      \begin{equation}
        E_\Pi[\Sigma]
        \equiv
        \int_\Sigma T^{(\Pi)}_{\mu\nu}\,\xi^\mu n^\nu\,\dd\Sigma,
        \label{eq:E-Pi-def}
      \end{equation}
      and its metric variation is
      \begin{equation}
        \delta E_\Pi
        = \half\int_\Sigma T^{(\Pi)}_{\mu\nu}\,\delta g^{\mu\nu}\,\dd\Sigma.
        \label{eq:delta-E-Pi}
      \end{equation}
      This definition is independent of $\beta(x)$ and of the heat-kernel structure.

    \subsection{Local spectral first law}
      \label{subsec:first-law}

      \begin{proposition}[Local spectral first law --- compatibility condition]
        \label{thm:first-law}
        In the infrared regime $R\ell_\chi^2 \ll 1$, and at leading order in the
        Seeley--DeWitt derivative expansion, the metric variation of the renormalized
        spectral entropy satisfies
        \begin{equation}
          \delta\Spi^{\mathrm ren}
          = \int_\Sigma \beta(x)^{-1}\,\delta E_\Pi(x)\,\dd^3 x
          \label{eq:first-law}
        \end{equation}
        as a structural compatibility identity between the geometric and matter sides
        of the variational problem.
        This identity holds if and only if the geometry satisfies the Einstein equation.
        It is therefore an equivalence condition, not an independent postulate.
      \end{proposition}

      \begin{proof}[Proof (compatibility argument)]
          From~\cite{Beau2026h}, the infrared-dominant local part of $\delta\Spi^{\mathrm ren}$
          contains the Einstein tensor:
          \begin{equation}
            \delta\Spi^{\mathrm ren}
            \supset
            \frac{1}{16\pi G_N}
            \int_M \dd^4x\,\sqrt{g}\;G_{\mu\nu}\,\delta g^{\mu\nu}
            + O(R^2\ell_\chi^2).
            \label{eq:deltaS-EH}
          \end{equation}
          The Bianchi identity $\nabla^\mu G_{\mu\nu} = 0$ implies that the integrand
          $G_{\mu\nu}\delta g^{\mu\nu}$ admits a local density representation on $\Sigma$
          once $\delta g^{\mu\nu}$ is supported in a compact region.
          On the matter side, from~\eqref{eq:delta-E-Pi}:
          \begin{equation}
            \delta E_\Pi
            = \half\int_\Sigma T^{(\Pi)}_{\mu\nu}\,\delta g^{\mu\nu}\,\dd\Sigma.
            \label{eq:delta-E-Pi-2}
          \end{equation}
          Requiring~\eqref{eq:first-law} to hold with $\beta(x)^{-1}$ as in~\eqref{eq:beta-decomp},
          then matching the two sides at leading IR order requires
          \begin{equation}
            \frac{1}{16\pi G_N}\,G_{\mu\nu}
            = \beta(x)^{-1} \cdot \half\,T^{(\Pi)}_{\mu\nu},
            \label{eq:matching}
          \end{equation}
          which is consistent with the Einstein equation when $\beta(x)^{-1}$ is evaluated
          at the equilibrium value.
          In spectral units where $\beta_0^{-1} = 1$, the curvature correction
          $-\frac{1}{6}R$ is of order $R\ell_\chi^2$ and is absorbed into the subleading
          remainder (see Remark~\ref{rem:gauge}).
          The identity~\eqref{eq:first-law} is therefore a structural compatibility condition:
          it holds if and only if the geometry satisfies the Einstein equation at leading
          infrared order.
          This establishes~\eqref{eq:first-law} as a compatibility condition rather than
          a postulate.
      \end{proof}

      \begin{remark}[Nature of the identity]
        \label{rem:identity-nature}
        Equation~\eqref{eq:first-law} is not a definition of $\beta(x)$ from $\delta\Spi$
        and $\delta E_\Pi$: both sides are independently defined.
        $\beta(x)^{-1}$ is defined from the heat kernel via Definition~\ref{def:beta};
        $\delta E_\Pi$ is defined from the matter action via~\eqref{eq:delta-E-Pi}.
        The identity~\eqref{eq:first-law} is therefore a non-trivial constraint, not a
        tautology.
      \end{remark}

  \section{Spectral equilibrium principle and Einstein equation}
    \label{sec:equilibrium}

    \subsection{Formulation of the equilibrium principle}
      \label{subsec:equilibrium-formulation}

      We now impose the following variational principle.

      \begin{definition}[Spectral equilibrium]
        \label{def:equilibrium}
        An admissible geometry $g$ is in \emph{spectral equilibrium} if it is a
        stationary point of the functional
        \begin{equation}
          \mathcal{F}[g,\psi]
          \equiv
          \Spi[g] - \beta_*^{-1}\,W_\Pi[g,\psi]
          \label{eq:F-functional}
        \end{equation}
        under arbitrary compactly supported metric variations $\delta g^{\mu\nu}$,
        with $\beta_*^{-1}$ a constant Lagrange multiplier.
      \end{definition}

      The functional $\mathcal{F}$ is the spectral analogue of a free energy: it extremizes
      entropic gain against the projected matter cost.
      The Lagrange multiplier $\beta_*^{-1}$ enforces the constraint that the projected
      matter content $W_\Pi$ is held fixed along admissible deformations.

      \begin{remark}[Nature of $\beta_*^{-1}$]
        \label{rem:beta-star}
        The global multiplier $\beta_*^{-1}$ in~\eqref{eq:F-functional} is distinct from
        the local field $\beta(x)^{-1}$ defined in Section~\ref{sec:spectral-locality}.
        The local field encodes spatial curvature variation; the global multiplier encodes
        the overall scale of the matter-geometry coupling.
        Their relation is discussed in Section~\ref{sec:normalization}.
      \end{remark}

    \subsection{Euler--Lagrange equation}
      \label{subsec:euler-lagrange}

      \begin{theorem}[Einstein equation as spectral equilibrium]
        \label{thm:einstein}
        In the infrared-dominant regime $R\ell_\chi^2 \ll 1$, the Euler--Lagrange
        equation of the spectral equilibrium condition $\delta\mathcal{F} = 0$ is
        \begin{equation}
          G_{\mu\nu} = 8\pi G_N\,T^{(\Pi)}_{\mu\nu},
          \label{eq:einstein}
        \end{equation}
        up to a cosmological term and corrections suppressed by $O(R\ell_\chi^2)$.
      \end{theorem}

      \begin{proof}
        Varying $\mathcal{F}[g,\psi]$ with respect to $g^{\mu\nu}$:
        \begin{equation}
          \delta\mathcal{F}
          = \delta\Spi - \beta_*^{-1}\,\delta W_\Pi = 0.
          \label{eq:delta-F}
        \end{equation}
        Substituting~\eqref{eq:deltaS-EH} and~\eqref{eq:W-variation}:
        \begin{equation}
          \frac{1}{16\pi G_N}\int_M\sqrt{g}\;G_{\mu\nu}\,\delta g^{\mu\nu}
          = \frac{\beta_*^{-1}}{2}\int_M\sqrt{g}\;T^{(\Pi)}_{\mu\nu}\,\delta g^{\mu\nu}
          + O(R\ell_\chi^2).
          \label{eq:EL-integrated}
        \end{equation}
        Since this holds for arbitrary compactly supported $\delta g^{\mu\nu}$, the
        fundamental lemma of the calculus of variations gives pointwise:
        \begin{equation}
          \frac{1}{16\pi G_N}\,G_{\mu\nu}
          = \frac{\beta_*^{-1}}{2}\,T^{(\Pi)}_{\mu\nu}
          + O(R\ell_\chi^2).
          \label{eq:EL-pointwise}
        \end{equation}
        For $\beta_*^{-1} = 1$ (see Section~\ref{sec:normalization}), this gives
        \begin{equation}
          G_{\mu\nu} = 8\pi G_N\,T^{(\Pi)}_{\mu\nu} + O(R\ell_\chi^2).
          \label{eq:einstein-proof}
        \end{equation}
      \end{proof}

      \begin{remark}[Absence of causal assumptions]
        \label{rem:no-horizon}
        The derivation requires no horizon structure, no local Rindler observers,
        no Raychaudhuri equation, and no assumption on causal structure.
        The field equation emerges from a purely variational condition on the spectral
        functional $\mathcal{F}$, defined on Riemannian manifolds.
        Its extension to Lorentzian signature is a separate question, discussed in
        Section~\ref{sec:discussion}.
      \end{remark}

  \section{Normalization of the factor \texorpdfstring{$8\pi G$}{8piG}}
    \label{sec:normalization}

    \subsection{Two independent normalizations}
      \label{subsec:two-norms}

      The factor $8\pi G$ in equation~\eqref{eq:einstein} is entirely determined by
      two independent normalization conventions:

      \paragraph{Geometric normalization.}
        The companion paper~\cite{Beau2026h} establishes that the infrared-dominant local
        part of the renormalized spectral entropy variation contains the Einstein tensor
        with coefficient $1/(16\pi G_N)$:
      \begin{equation}
        \frac{\delta\Spi^{\mathrm ren}}{\delta g^{\mu\nu}}\bigg|_{\mathrm IR,\,local}
        = \frac{1}{16\pi G_N}\,G_{\mu\nu}.
        \label{eq:geom-norm}
        \end{equation}
        This coefficient arises from the $a_2$ pole subtraction in the zeta-regularized
        determinant and the induced-gravity identification $G_N \sim 16\pi^2\ell_\chi^2$.

      \paragraph{Matter normalization.}
        The projected stress tensor is defined via the standard variational identity
      \begin{equation}
        T^{(\Pi)}_{\mu\nu} = \frac{2}{\sqrt{g}}
        \frac{\delta W_\Pi}{\delta g^{\mu\nu}},
        \label{eq:matter-norm}
        \end{equation}
        which carries the universal factor $2$ by convention.
        Equivalently, $\delta W_\Pi = \half\int\sqrt{g}\,T^{(\Pi)}_{\mu\nu}\delta g^{\mu\nu}$.

  \subsection{The factor $8\pi G$ as a product of normalizations}
    \label{subsec:8piG}

    \begin{proposition}[Normalization identity]
      \label{prop:normalization}
      The equilibrium condition $\delta\mathcal{F} = 0$ with $\beta_*^{-1} = 1$ gives
      \begin{equation}
        \frac{1}{16\pi G_N}\,G_{\mu\nu}
        = \half\,T^{(\Pi)}_{\mu\nu},
      \end{equation}
      which is equivalent to $G_{\mu\nu} = 8\pi G_N\,T^{(\Pi)}_{\mu\nu}$.
      The numerical factor $8\pi G_N$ arises as the unique product of two independent
      normalizations:
      \begin{equation}
        8\pi G_N
        = \frac{1}{(16\pi G_N)^{-1} \times \tfrac{1}{2}}
        = \frac{1}{(16\pi G_N)^{-1}} \cdot 2
        \cdot 1,
        \label{eq:8piG-factorization}
      \end{equation}
      where the factor $(16\pi G_N)^{-1}$ comes from the geometric normalization
      \eqref{eq:geom-norm} and the factor $\frac{1}{2}$ from the stress-tensor
      definition \eqref{eq:matter-norm}.
      More transparently: from $\frac{1}{16\pi G_N}G_{\mu\nu} = \frac{1}{2}T_{\mu\nu}$
      one reads off directly $G_{\mu\nu} = 8\pi G_N T_{\mu\nu}$.
      No other factor enters.
    \end{proposition}

  \subsection{Status of $\beta_*^{-1} = 1$}
    \label{subsec:beta-star-status}

    Setting $\beta_*^{-1} = 1$ in the constrained functional~\eqref{eq:F-functional}
    is not a dynamical assumption.
    It is a normalization choice equivalent to measuring energy in units where the
    spectral entropy and the projected matter functional are coupled with unit strength.
    Concretely, it amounts to absorbing $\beta_*^{-1}$ into the definition of
    $T^{(\Pi)}_{\mu\nu}$ as the effective stress tensor of the admissible projected
    sector, so that $W_\Pi$ is defined with the appropriate normalization.

    Alternatively, one may retain $\beta_*^{-1}$ as a free parameter and interpret
    $\beta_*^{-1} \neq 1$ as a departure from canonical spectral equilibrium.
    In that case,~\eqref{eq:EL-pointwise} generalizes to
    $G_{\mu\nu} = 8\pi G_N\,\beta_*^{-1}\,T^{(\Pi)}_{\mu\nu}$,
    which describes an off-equilibrium spectral geometry with an effective renormalization
    of the matter coupling.

  \subsection{Relation between $\beta_*^{-1}$ and the local field $\beta(x)^{-1}$}
    \label{subsec:beta-relation}

    The local field $\beta(x)^{-1} = \beta_0^{-1} - \frac{1}{6}R(x)$
    and the global Lagrange multiplier $\beta_*^{-1}$ play distinct roles.
    The local field encodes spatially varying curvature corrections to the spectral
    multiplier; the global multiplier controls the overall strength of the
    entropy-matter coupling in the equilibrium functional.

    Their connection arises when the equilibrium condition is imposed locally rather
    than globally.
    Requiring $\delta\mathcal{F}$ to vanish not only for global metric variations but
    also for local ones supported near a point $x$ replaces $\beta_*^{-1}$ by
    $\beta(x)^{-1}$ in~\eqref{eq:EL-pointwise}, yielding
    \begin{equation}
      G_{\mu\nu}(x)
      = 8\pi G_N\,\beta(x)^{-1}\,T^{(\Pi)}_{\mu\nu}(x) + O(R\ell_\chi^2).
      \label{eq:local-einstein}
    \end{equation}
    To display the structure of the curvature correction transparently, we
    choose \emph{spectral units} in which $\beta_0^{-1} = 1$.

    \begin{remark}[Gauge choice $\beta_0^{-1} = 1$ and the spectral scale $t_*$]
      \label{rem:gauge}
      Setting $\beta_0^{-1} = 1$ corresponds, via $\beta_0^{-1} = 2/t_*$, to the
      choice $t_* = 2$ in spectral units.
      This is a \emph{gauge choice for the admissibility scale}: it fixes the unit
      of spectral resolution but does not affect any physical observable.
      In particular, the induced Newton coupling $G_N \sim 16\pi^2 \ell_\chi^2$
      established in~\cite{Beau2026h} is independent of $t_*$, since it arises from
      the $a_2$ pole residue, which is a UV quantity fixed by the renormalization scheme.

      In these units, the local field simplifies to
      \begin{equation}
        \beta(x)^{-1} = 1 - \frac{1}{6}R(x) + O(\nabla^2 R),
        \label{eq:beta-spectral-units}
      \end{equation}
      and the locally imposed equilibrium condition~\eqref{eq:local-einstein} becomes
      \begin{equation}
        G_{\mu\nu}(x)
        = 8\pi G_N\!\left(1 - \tfrac{1}{6}R(x)\right) T^{(\Pi)}_{\mu\nu}(x)
        + O(R\ell_\chi^2).
        \label{eq:local-einstein-units}
      \end{equation}
      At zeroth order in $R\ell_\chi^2$, this reduces to the standard Einstein equation.
      The correction $-\frac{1}{6}R(x)$ is of order $R\ell_\chi^2$ and is absorbed into
      the $O(R\ell_\chi^2)$ remainder in the weak-curvature regime.
      Changing the gauge (i.e.\ changing $t_*$) rescales $\beta_0^{-1}$ and the
      correction simultaneously, leaving the physical content of~\eqref{eq:local-einstein}
      unchanged.
    \end{remark}

  \section{Relation to previous thermodynamic derivations}
    \label{sec:relation-jacobson}

    \subsection{Jacobson's derivation}
      \label{subsec:jacobson}

      Jacobson~\cite{Jacobson1995} derived Einstein's equation by applying the
      Clausius relation $\delta S = \delta Q / T$ to all local Rindler horizons.
      The key inputs are:
      (i) entropy proportional to horizon area, $S = A/(4G)$;
      (ii) temperature given by the Unruh formula, $T = a/(2\pi)$;
      (iii) heat flow $\delta Q = \int T_{\mu\nu}\chi^\mu\dd\Sigma^\nu$ across the horizon;
      (iv) the Raychaudhuri equation linking area change to matter flux.

      The present derivation parallels Jacobson's logical structure --- entropy, temperature,
      energy, equilibrium condition --- but replaces each ingredient with a spectral
      counterpart:
      (i) entropy is $\Spi = \frac{1}{2}\log\det' A_g$ (not an area functional);
      (ii) temperature is the local spectral multiplier $\beta(x)^{-1}$ (not Rindler);
      (iii) energy is the projected matter content $E_\Pi$ (not a horizon heat flux);
      (iv) equilibrium is a variational extremum (not a horizon balance condition).

      As a consequence, the present derivation is independent of the causal structure of
      spacetime: it holds on Riemannian manifolds and does not require a Lorentzian
      horizon.
      The extension to Lorentzian signature is discussed in Section~\ref{sec:discussion}.

    \subsection{Analogy with local Unruh temperature}
      \label{subsec:unruh-analogy}

      In the Unruh effect, the local temperature experienced by a Rindler observer with
      proper acceleration $a$ is $T_{\mathrm Unruh} = a/(2\pi)$.
      This is kinematic: it depends on the observer's trajectory.

      The local multiplier field $\beta(x)^{-1} = \beta_0^{-1} - \frac{1}{6}R(x)$
      has a formally analogous structure: a universal part $\beta_0^{-1}$ plus a local
      correction.
      However, the analogy is purely formal.
      The correction $-\frac{1}{6}R(x)$ is geometric (curvature-driven), not kinematic
      (acceleration-driven).
      No horizon is present or implied.
      The analogy should not be taken to mean that $\beta(x)^{-1}$ is a physical
      temperature in the Unruh sense: it is a spectral multiplier defined from
      heat-kernel data, not from a quantum state on a wedge.

    \subsection{Wald's Noether charge entropy}
      \label{subsec:wald}

      Wald~\cite{Wald1993} showed that black hole entropy for higher-derivative theories
      of gravity is given by a Noether charge, not by the area.
      The present framework is consistent with Wald's result in the following sense:
      the spectral entropy $\Spi$ is defined from the full operator $\Ag$, which encodes
      higher-derivative geometric information through the $a_4$ and higher Seeley--DeWitt
      coefficients.
      In the UV regime (or for $\xi \neq 0$), the dominant contribution to $\Spi$ shifts
      from the $a_2$-derived Einstein term to the $a_4$-derived quadratic curvature terms,
      whose variation would give a higher-derivative stress tensor consistent with Wald's
      formula.
      A detailed comparison lies beyond the scope of the present paper.

  \section{Discussion and outlook}
    \label{sec:discussion}

    \subsection{Summary}
      \label{subsec:summary}

      We have shown that the projective spectral entropy functional $\Spi[g]$ admits
      a natural local thermodynamic structure.
      The diagonal heat kernel defines a local spectral energy density $u(x;t)$ whose
      Seeley--DeWitt expansion yields a local multiplier field
      $\beta(x)^{-1} = \beta_0^{-1} - \frac{1}{6}R(x) + O(\nabla^2 R)$
      decomposed into a universal UV part and a geometric curvature correction.

      The local spectral first law $\delta\Spi = \int\beta^{-1}\delta E_\Pi$ is
      established as a compatibility theorem, not a definition: both sides are
      independently defined from spectral and matter data respectively.

      The spectral equilibrium principle --- extremizing $\Spi$ at fixed projected
      matter content --- yields Einstein's equation $G_{\mu\nu} = 8\pi G_N T_{\mu\nu}$
      as an Euler--Lagrange condition.
      The factor $8\pi G_N$ is fixed by the product of two universal normalizations from
      the geometric and matter sides, with no free parameter.

    \subsection{Limitations}
      \label{subsec:limitations}

      \paragraph{Euclidean setting.}
        The entire construction is performed on Riemannian manifolds.
        The extension to Lorentzian signature requires analytic continuation of the heat
        kernel, control over the retarded Green function, and a definition of $u(x;t)$ in
        the hyperbolic setting.
        This is the content of the Lorentzian completion program described in~\cite{Beau2026h}.

      \paragraph{Static foliation.}
        The projected matter energy $E_\Pi[\Sigma]$ is defined on a fixed spatial slice.
        A fully covariant formulation would require a foliation-independent definition,
        likely via the covariant phase space formalism.

      \paragraph{Higher Seeley--DeWitt corrections.}
        The equilibrium condition~\eqref{eq:einstein} holds at leading IR order.
        In the regime $R\ell_\chi^2 \sim 1$, higher Seeley--DeWitt contributions (from $a_4$
        and beyond) become comparable, and the field equation receives quadratic curvature
        corrections of the form $G_{\mu\nu} + \ell_\chi^2 B_{\mu\nu} + \cdots =
  8\pi G_N T_{\mu\nu}$.

      \paragraph{Matter coupling.}
        The projected stress tensor $T^{(\Pi)}_{\mu\nu}$ is defined formally from $W_\Pi$.
        An explicit construction of $W_\Pi$ in terms of standard matter fields requires
        identifying the admissible projected sector for each matter species.
        This is addressed in the Cosmochrony framework~\cite{Cosmochrony} but not derived
        here.

  \subsection{Outlook}
    \label{subsec:outlook}

    The present framework opens several directions.
    First, the Lorentzian completion of Paper~3 will test whether the equilibrium
    condition survives analytic continuation and whether it reproduces linearized
    gravitational wave propagation.
    Second, the higher Seeley--DeWitt regime ($a_4$ dominant, $\xi = 1/6$) may
    provide a spectral derivation of higher-curvature gravity theories.
    Third, the connection between the local field $\beta(x)^{-1}$ and the local
    Unruh effect could be made precise if a Lorentzian horizon is introduced as a
    special case of the general framework.

  \bibliographystyle{plain}
  \bibliography{cosmochrony-bibliography,external-refs}

\end{document}
